% ID: SEL-D0020-OA-E02
% ESTADO: validada. S021, revisada con los criterios oficiales.
\textbf{a)} La producción se restringe a $[0,2]$.
La derivada es $f'(x)=-3(x-1)(x-3)$. En $(0,1)$ ambos factores
entre paréntesis son negativos y la derivada es negativa. En $(1,2)$
tienen signos opuestos y la derivada es positiva. Así, el coste
decrece en $[0,1]$ y crece en $[1,2]$.

\textbf{b)} La derivada cambia de negativa a positiva en $x=1$.
El mínimo absoluto se alcanza al producir $\boxed{1\text{ tonelada}}$
y vale $f(1)=\boxed{26\text{ miles de euros}}$.

\textbf{c)} Por la monotonía, el máximo se alcanza en uno de los
extremos. Como $f(0)=30$ y $f(2)=28$, la producción que maximiza
el coste es $\boxed{0\text{ toneladas}}$.
